\section{Classical facts and normalizations}\label{sec:classical}

\subsection{Analytic low-rank inputs and parity}

We use modularity of elliptic curves over \(\mathbb Q\)~\cite{BCDT01}
and the Gross--Zagier--Kolyvagin theorem: analytic order zero or one
gives the corresponding rational rank and finiteness of the whole
Tate--Shafarevich group~\cite{GZ86,Kolyvagin90}. This implication
will be applied after the analytic order has been established.

Two quadratic-twist nonvanishing results supply the auxiliary analytic
inputs. For a primitive weight-two form of trivial central character,
Friedberg--Hoffstein supplies central values with prescribed local
behavior and compatible sign~\cite{FH95}; see also
\cite[Theorem~10.6]{AnandavardhananPrasad13}. In particular, when the
base form has odd functional sign, we may choose an odd fundamental
imaginary discriminant, split at the level and at any indicated finite
additional primes, whose twist has nonzero central value. When the base
central value is nonzero, the Bump--Friedberg--Hoffstein and
Murty--Murty derivative nonvanishing results supply imaginary Heegner
twists with a simple zero~\cite{BFH90,MurtyMurty91}.

The auxiliary fields may be chosen to have only the units \(\pm1\).
For a fundamental discriminant \(b\) coprime to the level \(N\),
the twist changes the functional sign by \(\chi_b(-N)\), where
\(\chi_b\) is the associated Dirichlet character. Every local
prescription below respects this sign condition.

Finally, the modular-symbol period theorem gives rationality of central
values after division by modular periods. Their nonnegativity for
elliptic newforms follows from the positivity theorem for self-dual
symplectic forms; see \cite{LapidRallis03}.

\begin{lemma}[Two-primary parity]\label{lem:classical-parity}
For every elliptic curve \(A/\mathbb Q\),
\[
 (-1)^{s_2(A)}=w(A).
\]
For a quadratic field \(K=\mathbb Q(\sqrt{k})\),
\[
 s_2(A/K)=s_2(A)+s_2(A^k).
\]
Moreover, \(\dim_{\mathbb F_2}\Sel_2(A)-\dim_{\mathbb F_2}A(\mathbb Q)[2]-s_2(A)\) is a nonnegative even integer. These statements do not require finiteness of \(\Sha(A)\).
\end{lemma}

\begin{proof}
Kramer's quadratic-extension parity comparison~\cite{Kramer81}
and the Kramer--Tunnell local root-number/norm-index
theorem~\cite{KramerTunnell82}, including the dyadic results
of \cite{DD10,DD11,CI16}, give by the Hilbert product formula
\[
 (-1)^{s_2(A)+s_2(A^b)}=w(A)w(A^b).
\]
Parity computed by finite two-descent agrees with Selmer-corank parity
after subtracting the rational two-torsion dimension. Indeed, the
Cassels--Tate pairing modulo divisibles is alternating, since the
elliptic polarization is represented by the rational divisor
\([O]\)~\cite[Remark~I.6.11]{MilneADT}
\cite[Theorem~5 and Corollary~7]{PoonenStoll1999}.

Calibrate the displayed relative parity equality using a coprime twist
of sign plus with nonzero central value, supplied by the nonvanishing
theorem. When the sign must change, take a negative discriminant split
at the level. The analytic rank-zero theorem gives \(s_2=0\) for this
twist, so its Selmer parity equals its functional sign. The relative
equality then proves the first assertion for \(A\).

Restriction and induction give the rational decomposition over a quadratic field. The associated isogeny preserves the rational Kummer Selmer conditions, proving the displayed corank identity. Finally the same alternating pairing gives the stated even difference between finite Selmer dimension and Selmer corank.
\end{proof}

\subsection{Gross--Zagier with absolute heights}

We use the conductor-one Gross--Zagier theorem at a coprime split
level, with the normalization in
\cite[Chapter~I, (4.3) and (6.4); Chapter~V, Theorem~(2.1)]{GZ86};
see also \cite{CST14,Conrad03}. The following choices fix the trace,
polarization, and height factors needed in the determinant comparison.

Let \(f\in S_2(\Gamma_0(N))\) be the normalized primitive eigenform,
and let \(K\) have odd fundamental discriminant \(-D<-4\), prime
to \(N\), with all level primes split. Choose an oriented ideal of
norm \(N\) with cyclic quotient. Let \(P_X\in J_0(N)(K)\) be the
Hilbert-class trace of the associated maximal-order Heegner point
minus the cusp \(\infty\); the trace sums once over the class group.
Write \((f,f)_N\) for the unnormalized Petersson integral on the
ordinary \(\Gamma_0(N)\)-quotient.

Project to the rational newform orbit orthogonally on the Jacobian,
up to isogeny, and use the \emph{restricted} polarization. Since
Rosati fixes the real Hecke field, its embeddings split the resulting
height pairing. Denote the diagonal at the embedding corresponding to
\(f\) by \(H_{\rm abs,f}\). Here absolute logarithmic heights are
divided by the field degree, and the pairing uses the Poincaré
principal-pairing convention.

\begin{theorem}[Split Gross--Zagier formula]\label{thm:classical-gz}
With the preceding hypotheses and units modulo signs one,
\[
 \frac{\sqrt D}{8\pi^2(f,f)_N}\,L'(f/K,1)=2H_{\rm abs,f}(P_X).
 \tag{GZ}\label{eq:GZ}
\]
\end{theorem}
The pairing in the relative formula is over \(K\), before division
by \([K:\mathbb Q]\); this accounts for the factor two in
\eqref{eq:GZ}. The ordinary pairing on the orbit projections is the
sum of the pairings at the Hecke embeddings.

For a rational modular parametrization \(\phi:X_0(N)\to E\), put
\(\phi^*\omega_E=c_E f\,dq/q\), where
\(q=\exp(2\pi i\tau)\), and let \(P=\phi_*P_X\). Then
\[
 L'(E/K,1)=
 \frac{\int_{E(\mathbb C)}|\omega_E\wedge\bar\omega_E|}
 {\sqrt D\,c_E^2}\cdot 2 H_{\rm abs}(P).
 \tag{GZ-E}\label{eq:GZ-E}
\]
Indeed, the projection formula divides the pushed-forward pairing by
\(\deg\phi\), while pullback of the differential gives
\[
 (\deg\phi)\int_{E(\mathbb C)}|\omega_E\wedge\bar\omega_E|
       =8\pi^2c_E^2(f,f)_N.
\]
These two degree factors cancel in \eqref{eq:GZ-E}.

The absolute Poincaré-pairing diagonal is
\(H_{\rm abs}(P)=\lim_j4^{-j}h_{x,\rm abs}(2^jP)\).
The origin divisor has canonical height \(H_{\rm abs}/2\), and
the Poincaré height uses its bilinearization \emph{without} another
halving. Thus this is the height convention fixed in the introduction.
The factor two from the relative height over \(K\) remains present.

We will compare \eqref{eq:GZ-E} with the BSD arithmetic volume using
the exact Manin, Tamagawa, and full-index factors for this single
Hilbert trace. This uses the height theorem; equality with the BSD
arithmetic volume is the further assertion to be proved.

\subsection{The split discrepancy identity}

Assume now that \(L(E/K,s)\) has a simple zero at one. The classical
elliptic theorem and isogeny give finiteness of the relevant
Tate--Shafarevich groups. Apply the abelian-variety isogeny invariance
of the BSD arithmetic volume~\cite[Theorem~2.1]{Tate1966}, extending
the elliptic comparison of Cassels~\cite{Cassels65}, to
\[
 \operatorname{Res}_{K/\mathbb Q}E_K\sim E\times E^k.
\]
We compute that volume on the restriction of scalars as follows.

The Néron restriction has period
\(\int|\omega_E\wedge\bar\omega_E|/\sqrt D\).
The differential remains minimal: bad places split, and good reduction
remains good. Coordinates in an integral quadratic basis give the
dual Lie-volume lattice; its real Jacobian contributes the quadratic
discriminant. Components at bad primes contribute \(c_\ell(E)^2\).
At ramified primes of good reduction, the special restriction fiber
is connected, with a vector kernel over the good special fiber.
Finally, the points and Tate--Shafarevich group are those over \(K\),
and the principal dual height diagonal is \(2H_{\rm abs}\).

\begin{corollary}[The full-index discrepancy]\label{cor:classical-G}
Under these assumptions, the single Hilbert-class trace satisfies
\[
  X(E)+X(E^k)
  =v_2\!\left(\frac{[E(K):\mathbb ZP]^2}
      {c_E^2\,\prod_\ell c_\ell(E)^2\ \#\Sha(E/K)}\right).       \tag{G}\label{eq:G}
\]
Here \([E(K):\mathbb ZP]\) is the full group index, including torsion.
\end{corollary}

\begin{proof}
Apply \eqref{eq:GZ-E}, then divide the analytic leading term by the arithmetic volume just computed. The squared full index converts the point height to the full Mordell--Weil regulator and retains the torsion factor. Taking \(v_2\) gives \eqref{eq:G}.
\end{proof}
The quantities defining \(X\) in the allowed analytic ranks are indeed positive rationals individually: in rank one take a nonvanishing split imaginary companion and use (GZ-E) and the rational elliptic rank-zero result for the other factor (or the modular period ratios). Here equality of arithmetic volumes on restriction and the product is over \(\mathbb R\), with \emph{total} real periods. The isogeny formula likewise allows changes of the rational curve in its isogeny class, or of its twists by induced isogenies, once in these analytic ranks. The general corank condition is itself isogeny-invariant without knowing these ranks. 
